% Plantilla para Trabajo Final (TF1 - 40%) del curso 1AEL0260
% Estructura completa segun syllabus: resumen, situacion problematica, fundamentos, causas,
% problemas de ingenieria, justificacion, estado del arte, arbol de objetivos,
% objetivos especificos, diagrama de bloques, cronograma Gantt, alcances, limitaciones,
% viabilidad, conclusiones.
% Uso: copiar a latex/proyecto/main.tex y rellenar
% Compilar: latexmk -pdf main.tex

\documentclass[12pt,a4paper]{article}

% --- Paquetes ---
\usepackage[utf8]{inputenc}
\usepackage[T1]{fontenc}
\usepackage[spanish]{babel}
\usepackage{graphicx}
\usepackage{amsmath}
\usepackage{booktabs}
\usepackage{hyperref}
\usepackage[margin=2.5cm]{geometry}
\usepackage{caption}
\usepackage{subcaption}
\usepackage{ganttchart}  % Para cronograma

% --- Metadatos ---
\title{[TITULO DEL PROYECTO DE TESIS]}
\author{Renzo Francisco Albatrino Aza \\
\texttt{renzo.albatrino@upc.edu.pe} \\
Universidad Peruana de Ciencias Aplicadas}
\date{UG-2do Semestre 2026}

\begin{document}

\maketitle

\begin{abstract}
[RESUMEN DEL PROYECTO - 1 parrafo con idea central, problema, propuesta de solucion y viabilidad]
\end{abstract}

% ============================================================
\section{Resumen del proyecto}
% ============================================================
[Breve resumen ejecutivo]

% ============================================================
\section{Situacion problematica de ingenieria}
% ============================================================
[Descripcion del problema de ingenieria en contexto academico y cientifico]

% ============================================================
\section{Fundamentos teoricos}
% ============================================================
[Conceptos teoricos que sustentan el proyecto]

% ============================================================
\section{Causas de la situacion problematica}
% ============================================================
[Analisis de causas que originan el problema]

% ============================================================
\section{Problemas de ingenieria}
% ============================================================
[Problemas especificos de ingenieria derivados]

% ============================================================
\section{Justificacion}
% ============================================================
[Por que es importante resolver este problema]

% ============================================================
\section{Estado del arte}
% ============================================================
\cite{[referencias del estado del arte]}
[Revision de trabajos previos, identificaccion de brecha]

% ============================================================
\section{Arbol de objetivos}
% ============================================================
[Descripcion del arbol de objetivos - problema general arriba, objetivos especificos abajo]

% ============================================================
\section{Objetivos especificos de desarrollo e implementacion}
% ============================================================

\subsection{Objetivo general}
[OBJETIVO GENERAL]

\subsection{Objetivos especificos}
\begin{itemize}
    \item \textbf{OE1:} [objetivo 1]
    \item \textbf{OE2:} [objetivo 2]
    \item \textbf{OE3:} [objetivo 3]
    \item \textbf{OE4:} [objetivo 4]
    \item \textbf{OE5:} [objetivo 5]
\end{itemize}

% ============================================================
\section{Diagrama de bloques}
% ============================================================
[Descripcion de la arquitectura por bloques]

% \begin{figure}[ht]
%     \centering
%     \includegraphics[width=\textwidth]{figuras/diagrama_bloques}
%     \caption{Diagrama de bloques de la solucion propuesta.}
%     \label{fig:bloques}
% \end{figure}

% ============================================================
\section{Cronograma por objetivo especifico}
% ============================================================

% \begin{ganttchart}[
%     y unit title=0.7cm,
%     bar/.append style={fill=blue!30},
%     bar height=0.5,
%     group label font=\small,
%     bar label font=\small
% ]{1}{14}
%     \gantttitle{Cronograma 2026} \\
%     \ganttbar[bar/.append style={fill=red!40}]{OE1: Estado del arte}{1}{3} \\
%     \ganttbar[bar/.append style={fill=green!40}]{OE2: Arquitectura}{4}{8} \\
%     \ganttbar[bar/.append style={fill=blue!40}]{OE3: Implementacion}{9}{13} \\
%     \ganttbar[bar/.append style={fill=orange!40}]{OE4: Evaluacion}{13}{15} \\
%     \ganttbar[bar/.append style={fill=purple!40}]{OE5: Viabilidad}{15}{16} \\
% \end{ganttchart}

% ============================================================
\section{Alcances y limitaciones}
% ============================================================

\subsection{Alcances}
[Que abarca el proyecto]

\subsection{Limitaciones}
[Que no abarca el proyecto - ser honesto]

% ============================================================
\section{Viabilidad}
% ============================================================

\subsection{Viabilidad tecnica}
[Analisis de viabilidad tecnica]

\subsection{Viabilidad economica}
[Analisis de viabilidad economica]

\subsection{Viabilidad social}
[Analisis de viabilidad social]

\subsection{Viabilidad operativa}
[Analisis de viabilidad operativa]

% ============================================================
\section{Conclusiones y comentarios finales}
% ============================================================
[Conclusiones del proyecto]

\bibliographystyle{ieeetr}
\bibliography{references}

\end{document}
